% ===============================================================
% GRADE 7 -- MATATAG Science 7, Three-Term BOW (rev. 8 Apr 2026)
% ===============================================================
\section{Grade 7}
\resetcards
\begin{gradebody}

% ---------------------------------------------------------------
\subsection{Term 1: Matter, Investigations, Solutions, Cells}

\topic{Models and the Particle Model of Matter}
\kf
\begin{keyfacts}
\item A \textbf{model} is a simplified representation of something too small, large, or hidden to observe directly.
\item \textbf{Particle model:} all matter is made of tiny particles, and each pure substance has its own kind.
\item Particles are always moving, have spaces between them, attract each other, and move faster when heated.
\item \textbf{Solid:} packed, fixed positions, vibrate only, so it has definite shape and volume.
\item \textbf{Liquid:} close but slide past each other, so it has definite volume and takes the container's shape.
\item \textbf{Gas:} far apart and moving fast and randomly, so it has no fixed shape or volume and is compressible.
\item Heat \emph{absorbed}: melting, evaporation/boiling, sublimation. Heat \emph{released}: freezing, condensation, deposition.
\item A pure substance stays at constant temperature while it changes state.
\end{keyfacts}
\cards
\qa{Simplified representation of something hard to observe}{Model}{E}
\qa{State with fixed shape and fixed volume}{Solid}{E}
\qa{State most easily compressed}{Gas}{E}
\qa{Particles only vibrate in fixed positions}{Solid}{A}
\qa{Solid changes directly to gas}{Sublimation}{E}
\qa{Gas changes directly to solid}{Deposition}{A}
\qa{Mothballs shrinking in a closet}{Sublimation}{A}
\qa{Droplets forming outside a glass of iced water}{Condensation}{E}
\qa{Effect of heating on particle motion}{Particles move faster}{E}
\qa{Food colouring spreads faster in hot or cold water?}{Hot water}{A}
\qa{Temperature of pure ice while it is melting}{Stays constant (\SI{0}{\celsius})}{D}
\qa{Why gases can be compressed}{Large spaces between particles}{A}
\qa{Changes of state that release heat}{Freezing, condensation, deposition}{D}
\begin{confbox}
\confusion{Evaporation}{boiling}{evaporation happens only at the surface and at any temperature; boiling happens throughout the liquid at the boiling point.}
\confusion{Melting}{dissolving}{melting is solid $\rightarrow$ liquid by heating one substance; dissolving is a solute spreading through a solvent.}
\confusion{Sublimation}{deposition}{sublimation is solid $\rightarrow$ gas; deposition is gas $\rightarrow$ solid.}
\end{confbox}

\topic{Scientific Investigations and Measurement}
\kf
\begin{keyfacts}
\item Report parts (BOW): aim/problem $\rightarrow$ materials and equipment $\rightarrow$ method $\rightarrow$ results (data) $\rightarrow$ conclusion.
\item \textbf{Independent} variable is what you change; \textbf{dependent} is what you measure; \textbf{controlled} variables are kept the same.
\item SI base units: m, kg, s, K, mol, A, cd (see Cheat Sheets).
\item Read liquid volume at the \emph{bottom of the meniscus}, with your eye level to it.
\item Repeat trials to make results more reliable.
\end{keyfacts}
\cards
\qa{Variable deliberately changed}{Independent variable}{E}
\qa{Variable measured as the result}{Dependent variable}{E}
\qa{Factors kept the same in a fair test}{Controlled variables}{E}
\qa{SI base unit of mass}{Kilogram}{E}
\qa{SI base unit of temperature}{Kelvin}{A}
\qa{SI base unit of amount of substance}{Mole}{D}
\qa{Tool for measuring liquid volume}{Graduated cylinder}{E}
\qa{Where to read a water meniscus}{Bottom of the curve, at eye level}{A}
\qa{\SI{1}{\liter} equals}{\SI{1000}{\milli\liter}}{E}
\qa{Prefix meaning one-thousandth}{milli-}{A}
\qa{Part of a report that answers the aim}{Conclusion}{E}
\qa{Testing how water temperature affects sugar's dissolving time: dependent variable?}{Time to dissolve}{D}
\qa{Repeating trials improves the}{Reliability}{A}
\begin{confbox}
\confusion{Accuracy}{precision}{accuracy means close to the true value; precision means repeated readings are close to each other.}
\confusion{Reliability}{validity}{reliable results repeat consistently; a valid test actually measures what it claims to.}
\end{confbox}

\topic{Solutions, Solubility, Acids and Bases}
\kf
\begin{keyfacts}
\item A \textbf{solution} is a homogeneous mixture of a \textbf{solute} (dissolved, usually the smaller amount) in a \textbf{solvent} (usually the larger amount).
\item Water is the ``universal solvent.''
\item Concentration: $\%\text{ by mass} = \dfrac{\text{mass of solute}}{\text{mass of solution}} \times 100$.
\item \textbf{Unsaturated} solutions can dissolve more solute; \textbf{saturated} ones cannot; \textbf{supersaturated} ones hold more than normal (unstable).
\item Solubility depends on temperature, the nature of solute and solvent (``like dissolves like''), and pressure (for gases).
\item Most solids dissolve better when hot, but gases dissolve \emph{less} when warm.
\item Stirring and crushing speed up \emph{dissolving}; they do not change solubility.
\item Litmus: acid turns blue litmus red; base turns red litmus blue; neutral salt solutions cause no change.
\end{keyfacts}
\cards
\qa{Component present in the greater amount}{Solvent}{E}
\qa{Universal solvent}{Water}{E}
\qa{Solution that can still dissolve more solute}{Unsaturated}{E}
\qa{Holds more solute than normally possible}{Supersaturated}{A}
\qa{Blue litmus dipped in vinegar turns}{Red}{E}
\qa{Red litmus dipped in soap solution turns}{Blue}{E}
\qa{pH of pure water at \SI{25}{\celsius}}{7 (neutral)}{A}
\qa{Why a warm soft drink goes flat faster}{Gas solubility falls as temperature rises}{D}
\qa{Stirring changes solubility or rate of dissolving?}{Rate of dissolving only}{D}
\qa{\SI{20}{g} sugar in \SI{80}{g} water: \% by mass}{20\%}{D}
\qa{Air is a gaseous solution; its solvent}{Nitrogen}{D}
\qa{Brass is a solid solution of}{Copper and zinc}{A}
\qa{Does oil dissolve in water?}{No (``like dissolves like'')}{A}
\begin{confbox}
\confusion{Solubility}{rate of dissolving}{solubility is \emph{how much} dissolves; rate is \emph{how fast}.}
\confusion{Concentrated}{saturated}{concentrated means a lot of solute; saturated means no more can dissolve at that temperature.}
\end{confbox}
\mnemonic{\textbf{BRA}: \textbf{B}lue $\rightarrow$ \textbf{R}ed = \textbf{A}cid (so red $\rightarrow$ blue = base).}

\topic{The Compound Microscope}
\kf
\begin{keyfacts}
\item Total magnification = eyepiece $\times$ objective (e.g., $10\times \cdot 40\times = 400\times$).
\item Objectives: scanner ($4\times$), LPO ($10\times$), HPO ($40\times$), and oil immersion ($100\times$, needs oil).
\item \textbf{Diaphragm} controls light; \textbf{stage clips} hold the slide; \textbf{nosepiece} rotates the objectives.
\item Coarse knob is for the lowest power only; use only the \textbf{fine} knob on HPO.
\item Carry with one hand on the arm and one under the base. Start with the lowest-power objective.
\item The image is inverted (upside down and reversed). Moving to higher power gives a smaller, dimmer field of view.
\end{keyfacts}
\cards
\qa{Magnification with $10\times$ eyepiece and $40\times$ objective}{$400\times$}{E}
\qa{Knob to use with the HPO}{Fine adjustment}{A}
\qa{Part that regulates light}{Diaphragm}{E}
\qa{Holds the slide in place}{Stage clips}{E}
\qa{Rotating part holding the objectives}{Revolving nosepiece}{E}
\qa{Lens nearest the eye}{Eyepiece (ocular)}{E}
\qa{Letter ``e'' as seen through the microscope}{Upside down and reversed}{A}
\qa{Correct way to carry it}{One hand on arm, one under base}{E}
\qa{Objective to focus with first}{Lowest power (scanner)}{A}
\qa{LPO to HPO: field of view becomes}{Smaller and darker}{D}
\qa{The $100\times$ objective requires}{Immersion oil}{D}
\begin{confbox}
\confusion{Coarse}{fine adjustment}{coarse moves the stage a lot (low power only); fine makes small focus changes.}
\confusion{Magnification}{resolution}{magnification enlarges the image; resolution separates two close points clearly.}
\end{confbox}

\topic{Plant and Animal Cells}
\kf
\begin{keyfacts}
\item The cell is the basic unit of life. \textbf{Robert Hooke} named ``cells'' from cork (1665).
\item Cell theory: \textbf{Schleiden} (plants), \textbf{Schwann} (animals), \textbf{Virchow} (cells come from pre-existing cells).
\item \textbf{Nucleus} is the control centre (DNA); \textbf{cell membrane} controls entry and exit; \textbf{cytoplasm} is the jelly-like interior.
\item \textbf{Mitochondria} carry out respiration (the ``powerhouse''); \textbf{chloroplasts} carry out photosynthesis; \textbf{ribosomes} make proteins.
\item Plant cells only: cell wall (cellulose), chloroplasts, and a large central vacuole. Animal cells are irregular and have small vacuoles.
\item Organisms are \textbf{unicellular} (bacteria, amoeba) or \textbf{multicellular} (humans).
\end{keyfacts}
\cards
\qa{Who coined the term ``cell''}{Robert Hooke}{E}
\qa{``All cells come from pre-existing cells''}{Rudolf Virchow}{A}
\qa{Proposed that all plants are made of cells}{Matthias Schleiden}{A}
\qa{Powerhouse of the cell}{Mitochondrion}{E}
\qa{Organelle for photosynthesis}{Chloroplast}{E}
\qa{Control centre of the cell}{Nucleus}{E}
\qa{Material of the plant cell wall}{Cellulose}{A}
\qa{Organelle that makes proteins}{Ribosome}{A}
\qa{Example of a unicellular organism}{Bacterium / amoeba}{E}
\qa{Cell without a true nucleus}{Prokaryotic cell}{D}
\qa{Plant organelle that keeps the cell firm (turgid)}{Large central vacuole}{D}
\qa{Organelle absent in onion bulb epidermis cells}{Chloroplast}{D}
\qa{Common stain for cheek cells}{Methylene blue}{A}
\begin{confbox}
\confusion{Cell wall}{cell membrane}{the wall is a rigid outer layer (plants only); the membrane is a flexible, selective barrier (all cells).}
\confusion{Chloroplast}{chlorophyll}{the chloroplast is the organelle; chlorophyll is the green pigment inside it.}
\confusion{Prokaryote}{eukaryote}{a prokaryote has no nucleus (bacteria); a eukaryote has a nucleus (plants, animals, fungi).}
\end{confbox}

% ---------------------------------------------------------------
\subsection{Term 2: Reproduction, Ecology, Forces, Motion}

\topic{Cellular Reproduction}
\kf
\begin{keyfacts}
\item \textbf{Mitosis} gives 2 identical cells and is used for growth and repair in body cells.
\item \textbf{Meiosis} gives 4 genetically different sex cells (gametes) with half the chromosomes.
\item Human body cell: 46 chromosomes; gamete: 23.
\item Stages: interphase (DNA copies), then prophase, metaphase, anaphase, telophase, then cytokinesis.
\item Fertilization: sperm + egg $\rightarrow$ zygote.
\item \textbf{Sexual} reproduction: 2 parents, varied offspring. \textbf{Asexual}: 1 parent, identical offspring.
\item Asexual types: binary fission (bacteria), budding (yeast, hydra), spores (ferns, fungi), vegetative propagation (cuttings), fragmentation.
\end{keyfacts}
\cards
\qa{Cell division for growth and repair}{Mitosis}{E}
\qa{Cell division that makes gametes}{Meiosis}{E}
\qa{Daughter cells from one meiosis}{Four}{E}
\qa{Chromosomes in a human sperm cell}{23}{A}
\qa{Cell formed when sperm and egg fuse}{Zygote}{E}
\qa{Chromosomes line up at the middle}{Metaphase}{A}
\qa{Sister chromatids pulled apart}{Anaphase}{A}
\qa{Stage when DNA is copied}{Interphase}{D}
\qa{Division of the cytoplasm}{Cytokinesis}{A}
\qa{Yeast reproduces asexually by}{Budding}{A}
\qa{Bacteria reproduce by}{Binary fission}{A}
\qa{Offspring genetically identical to the parent}{Asexual reproduction}{E}
\qa{Ferns and mushrooms reproduce asexually by}{Spores}{A}
\qa{Why siblings are not identical}{Meiosis + random fertilization}{D}
\begin{confbox}
\confusion{Mitosis}{meiosis}{mitosis: 1 division, 2 identical body cells. Meiosis: 2 divisions, 4 varied gametes.}
\confusion{Diploid}{haploid}{diploid cells have 2 sets (46 in humans); haploid gametes have 1 set (23).}
\confusion{Chromosome}{chromatid}{a chromatid is one of the two identical copies of a replicated chromosome, joined at the centromere.}
\end{confbox}
\mnemonic{\textbf{IPMAT}: Interphase, Prophase, Metaphase (\textbf{M}iddle), Anaphase (\textbf{A}part), Telophase (\textbf{T}wo nuclei).}

\topic{Levels of Biological Organization}
\kf
\begin{keyfacts}
\item Cell $\rightarrow$ tissue $\rightarrow$ organ $\rightarrow$ organ system $\rightarrow$ organism $\rightarrow$ population $\rightarrow$ community $\rightarrow$ ecosystem $\rightarrow$ (biome) $\rightarrow$ biosphere.
\item \textbf{Population}: one species in one area. \textbf{Community}: all the populations that interact there.
\item \textbf{Ecosystem} = community + nonliving (abiotic) factors. \textbf{Biosphere}: all life on Earth.
\end{keyfacts}
\cards
\qa{Group of similar cells doing one job}{Tissue}{E}
\qa{Same species living in one area}{Population}{E}
\qa{Community plus nonliving factors}{Ecosystem}{E}
\qa{Highest level of organization}{Biosphere}{E}
\qa{The heart is an example of}{An organ}{E}
\qa{Level right after organ}{Organ system}{E}
\qa{All tamaraws in Mindoro}{Population}{A}
\qa{Fish, algae, corals, water, and sunlight in a reef}{Ecosystem}{A}
\qa{Blood is classified as a}{Tissue (connective)}{D}
\begin{confbox}
\confusion{Population}{community}{a population is one species; a community is many species living together.}
\confusion{Community}{ecosystem}{an ecosystem adds the abiotic factors (water, soil, light).}
\end{confbox}

\topic{Trophic Levels and Energy Transfer}
\kf
\begin{keyfacts}
\item The Sun is the ultimate energy source. Producers (autotrophs) form the 1st trophic level.
\item Primary consumers are herbivores; secondary consumers eat herbivores; then tertiary consumers.
\item Decomposers (bacteria, fungi) recycle dead matter.
\item \textbf{10\% rule:} only about 10\% of energy passes to the next level. The rest is used or lost as heat.
\item Energy loss limits food chains to about 4--5 levels. A pyramid of energy is always upright.
\end{keyfacts}
\cards
\qa{Organisms at the 1st trophic level}{Producers}{E}
\qa{Grass $\rightarrow$ grasshopper $\rightarrow$ frog $\rightarrow$ snake: secondary consumer}{Frog}{A}
\qa{Energy passed to the next level}{About 10\%}{E}
\qa{Form of most energy ``lost'' between levels}{Heat}{A}
\qa{Organisms that break down dead matter}{Decomposers}{E}
\qa{Another name for primary consumers}{Herbivores}{E}
\qa{Interconnected food chains}{Food web}{E}
\qa{Producers hold \SI{10000}{kJ}; primary consumers get about}{\SI{1000}{kJ}}{D}
\qa{Why food chains rarely exceed 5 levels}{Too little energy left}{D}
\qa{Scientist linked to the 10\% rule}{Raymond Lindeman}{D}
\begin{confbox}
\confusion{Food chain}{food web}{a chain is one feeding path; a web is many interlinked chains.}
\confusion{Pyramid of energy}{pyramid of numbers}{the energy pyramid is always upright; a numbers pyramid can be inverted (one tree, many insects).}
\end{confbox}

\topic{Balanced and Unbalanced Forces}
\kf
\begin{keyfacts}
\item A force is a push or pull, measured in \textbf{newtons (N)} with a spring balance.
\item \textbf{Balanced} forces (net force 0) mean the object is at rest or moving at constant velocity.
\item \textbf{Unbalanced} forces change speed or direction (acceleration).
\item A free-body diagram uses arrows whose length shows size and whose head shows direction.
\item Weight $W = mg$, with $g \approx \SI{9.8}{\meter\per\second\squared}$.
\end{keyfacts}
\cards
\qa{SI unit of force}{Newton}{E}
\qa{Device for measuring force}{Spring balance}{E}
\qa{Net force on a car at constant velocity}{Zero}{A}
\qa{Freely falling mango: balanced or unbalanced?}{Unbalanced}{E}
\qa{Box resting on a ramp: balanced or unbalanced?}{Balanced}{A}
\qa{Diagram of all forces on one object}{Free-body diagram}{E}
\qa{Force that opposes motion between surfaces}{Friction}{E}
\qa{Support force perpendicular to a surface}{Normal force}{A}
\qa{\SI{10}{N} right and \SI{4}{N} left: net force}{\SI{6}{N} to the right}{A}
\qa{Unbalanced forces change an object's}{Speed or direction}{A}
\qa{Weight of a \SI{2}{kg} object ($g = 9.8$)}{\SI{19.6}{N}}{D}
\begin{confbox}
\confusion{Mass}{weight}{mass is the amount of matter (kg, same everywhere); weight is the pull of gravity (N, changes with $g$).}
\confusion{Balanced forces}{no forces}{a book at rest still has forces acting on it; they just cancel.}
\end{confbox}

\topic{Motion: Distance, Displacement, Speed, Velocity}
\kf
\begin{keyfacts}
\item Position is described from a \textbf{reference point}.
\item \textbf{Distance} is total path length (scalar). \textbf{Displacement} is the straight-line change in position with direction (vector).
\item Speed $v = d/t$ (scalar). Velocity = displacement/time (vector).
\item On a distance--time graph, a straight sloped line means uniform velocity, the slope is the speed, and a flat line means at rest.
\item $\SI{1}{\meter\per\second} = \SI{3.6}{km/h}$.
\end{keyfacts}
\cards
\qa{Quantity with magnitude and direction}{Vector}{E}
\qa{Is speed scalar or vector?}{Scalar}{E}
\qa{Point used to describe position}{Reference point}{E}
\qa{\SI{3}{m} east then \SI{3}{m} west: distance, displacement}{\SI{6}{m}; \SI{0}{m}}{A}
\qa{Car covers \SI{100}{m} in \SI{5}{s}: speed}{\SI{20}{\meter\per\second}}{E}
\qa{Slope of a distance--time graph}{Speed}{A}
\qa{Flat (horizontal) line on a d--t graph}{At rest}{A}
\qa{Straight sloped line on a d--t graph}{Uniform velocity}{A}
\qa{Steeper d--t line means}{Faster}{A}
\qa{\SI{72}{km/h} in m/s}{\SI{20}{\meter\per\second}}{D}
\begin{confbox}
\confusion{Distance}{displacement}{distance is the whole path; displacement is start to finish with direction.}
\confusion{Speed}{velocity}{velocity includes direction, so a turning car at constant speed has changing velocity.}
\end{confbox}

% ---------------------------------------------------------------
\subsection{Term 3: Heat, Earthquakes, the Sun and Weather}

\topic{Heat Transfer}
\kf
\begin{keyfacts}
\item \textbf{Heat} is energy transferred because of a temperature difference (joules). \textbf{Temperature} measures the average kinetic energy of particles (\si{\celsius}, K).
\item Heat flows from hot to cold.
\item \textbf{Conduction} is transfer by particle collisions in contact (best in metals).
\item \textbf{Convection} is transfer by moving fluids: warm fluid is less dense and rises.
\item \textbf{Radiation} is transfer by electromagnetic waves and needs no medium (Sun to Earth).
\item $K = {}^\circ\text{C} + 273$.
\item Heat-to-electricity devices: thermoelectric generators and geothermal power plants (e.g., Tiwi, Albay).
\end{keyfacts}
\cards
\qa{Metal spoon heating up in hot soup}{Conduction}{E}
\qa{How the Sun's heat reaches Earth}{Radiation}{E}
\qa{Circulating currents in boiling water}{Convection}{E}
\qa{Transfer that works in a vacuum}{Radiation}{A}
\qa{Why pot handles are wood or plastic}{They are insulators}{E}
\qa{SI unit of heat}{Joule}{E}
\qa{Measure of average particle kinetic energy}{Temperature}{A}
\qa{\SI{25}{\celsius} in kelvin}{\SI{298}{K}}{A}
\qa{Why warm fluids rise}{Less dense}{A}
\qa{Dark shirts get hotter in sunlight because they}{Absorb more radiation}{A}
\qa{Why air conditioners are mounted high}{Cool air sinks (convection)}{D}
\qa{Device turning a temperature difference directly into electricity}{Thermoelectric generator}{D}
\qa{Geothermal plant in Albay}{Tiwi}{D}
\begin{confbox}
\confusion{Heat}{temperature}{heat is energy in transit (J); temperature is how hot (\si{\celsius}, K).}
\confusion{Conduction}{convection}{conduction is particles passing energy without flowing; convection is the fluid itself moving.}
\end{confbox}

\topic{Faults, Earthquakes, and Tsunamis}
\kf
\begin{keyfacts}
\item A \textbf{fault} is a crack in the crust along which blocks have moved. The block above an inclined fault is the hanging wall; the block below is the footwall.
\item \textbf{Normal} fault: hanging wall moves down (tension). \textbf{Reverse/thrust}: hanging wall moves up (compression). \textbf{Strike-slip}: horizontal sliding (shear).
\item \textbf{Focus} (hypocenter) is where the rupture starts underground; the \textbf{epicenter} is directly above it on the surface.
\item \textbf{Magnitude} is the energy released (one number per quake). \textbf{Intensity} is the shaking and damage at a place; PHIVOLCS uses PEIS I--X.
\item The Philippine Fault Zone is a left-lateral strike-slip fault about \SI{1200}{km} long. The West Valley Fault threatens Metro Manila.
\item Find the nearest active fault with \textbf{PHIVOLCS FaultFinder}.
\item Tsunamis are mostly caused by undersea quakes that move the seafloor vertically. A natural warning is the sea suddenly pulling back.
\item During shaking: \textbf{Duck, Cover, and Hold}.
\end{keyfacts}
\cards
\qa{Point on the surface above the focus}{Epicenter}{E}
\qa{Underground origin of a quake}{Focus (hypocenter)}{E}
\qa{Fault where the hanging wall drops}{Normal fault}{E}
\qa{Fault where the hanging wall rises}{Reverse fault}{E}
\qa{Blocks slide horizontally}{Strike-slip fault}{E}
\qa{Stress that forms normal faults}{Tension}{A}
\qa{Stress that forms reverse faults}{Compression}{A}
\qa{Movement type of the Philippine Fault Zone}{Left-lateral strike-slip}{D}
\qa{Approximate length of the Philippine Fault Zone}{About \SI{1200}{km}}{D}
\qa{PHIVOLCS intensity scale}{PEIS (I to X)}{A}
\qa{Magnitude measures}{Energy released}{A}
\qa{Instrument recording ground shaking}{Seismograph}{E}
\qa{PH agency for quakes and volcanoes}{PHIVOLCS}{E}
\qa{What to do during shaking}{Duck, Cover, and Hold}{E}
\qa{Natural sign of an incoming tsunami}{Sea suddenly recedes}{A}
\qa{Fault system threatening Metro Manila}{West Valley Fault}{A}
\qa{PHIVOLCS tool to find the nearest fault}{FaultFinder}{A}
\qa{Fault type most likely to cause a tsunami}{Reverse/thrust (vertical seafloor shift)}{D}
\qa{Magnitude of the 2013 Bohol earthquake}{7.2}{D}
\qa{Fault types most common in PH, and why}{\verify{Reverse/thrust and strike-slip, from plate convergence (expected DepEd answer)}}{D}
\begin{confbox}
\confusion{Magnitude}{intensity}{magnitude is one value for the whole quake; intensity differs from place to place.}
\confusion{Focus}{epicenter}{the focus is underground; the epicenter is on the surface above it.}
\confusion{Tsunami}{storm surge}{a tsunami comes from seafloor displacement (quake or eruption); a storm surge is seawater pushed by typhoon winds and low pressure.}
\end{confbox}

\topic{The Sun's Influence: Seasons, Breezes, Monsoons, ITCZ}
\kf
\begin{keyfacts}
\item Earth's axis is tilted about \SI{23.5}{\degree}. Rotation (about \SI{24}{h}) causes day and night. Revolution (about $365\frac{1}{4}$ days) plus the tilt causes the seasons.
\item Direct rays are more intense; slanting rays spread over a larger area.
\item The PH is in the Northern Hemisphere: longest day at the June solstice, shortest at the December solstice. Equinoxes (March and September) have about equal day and night.
\item Land heats and cools faster than water. Daytime \textbf{sea breeze} blows sea $\rightarrow$ land; nighttime \textbf{land breeze} blows land $\rightarrow$ sea.
\item \textbf{Amihan} (northeast monsoon) is cool and dry. \textbf{Habagat} (southwest monsoon) is warm, moist, and rainy.
\item \textbf{ITCZ}: a low-pressure belt near the equator where the trade winds meet, bringing clouds and rain.
\item The ozone layer absorbs most UV radiation.
\end{keyfacts}
\cards
\qa{Main cause of the seasons}{Tilt of Earth's axis}{E}
\qa{Tilt of Earth's axis}{About \SI{23.5}{\degree}}{E}
\qa{Day and night are caused by}{Rotation}{E}
\qa{Time for one revolution}{About $365\frac{1}{4}$ days}{E}
\qa{Daytime breeze from sea to land}{Sea breeze}{E}
\qa{Why land heats faster than the sea}{Water has a higher heat capacity}{A}
\qa{Local name of the northeast monsoon}{Amihan}{E}
\qa{Local name of the southwest monsoon}{Habagat}{E}
\qa{Monsoon behind heavy PH rains in the wet season}{Habagat}{A}
\qa{Belt where the trade winds converge}{ITCZ}{A}
\qa{Longest daytime in the PH}{June solstice}{A}
\qa{About equal day and night}{Equinox}{E}
\qa{Gas layer that absorbs UV}{Ozone}{A}
\qa{Why slanting rays heat less}{Energy spread over a larger area}{D}
\qa{Earth is closest to the Sun in January, yet it is the cool season in the PH. This shows}{Seasons aren't caused by distance}{D}
\begin{confbox}
\confusion{Rotation}{revolution}{rotation is spinning on the axis (day/night); revolution is orbiting the Sun (year, seasons).}
\confusion{Sea breeze}{land breeze}{a breeze is named after where it comes \emph{from}: sea breeze by day, land breeze by night.}
\confusion{Amihan}{Habagat}{Amihan is NE, cool and dry, about Nov--Feb; Habagat is SW, rainy, about Jun--Sep.}
\confusion{Weather}{climate}{weather is short-term conditions; climate is the long-term (about 30-year) average.}
\end{confbox}

\end{gradebody}
\cardstats
